\documentclass[12pt,a4paper]{article}

%!TEX program = pdflatex

% =============================================================================
% 1. PHONG CHU & NGON NGU
% =============================================================================
\usepackage[utf8]{inputenc}
\usepackage[T5]{fontenc}
\usepackage[vietnamese]{babel}

% =============================================================================
% 2. BO CUC TRANG
% =============================================================================
\usepackage[margin=2.5cm]{geometry}
\usepackage{setspace}
\onehalfspacing

% =============================================================================
% 3. BANG BIEU & DANH SACH
% =============================================================================
\usepackage{array}
\usepackage{booktabs}
\usepackage{longtable}
\usepackage{enumitem}
\usepackage{tabularx}
\usepackage{makecell}
\usepackage{float}
\usepackage{etoolbox}
\usepackage{amssymb}

% =============================================================================
% 4. TIEN ICH VAN BAN
% =============================================================================
\usepackage{xcolor}
\usepackage{titlesec}
\usepackage{tocloft}
\usepackage{xurl}
\usepackage{seqsplit}
\usepackage{hyperref}

\hypersetup{
    colorlinks=false,
    hidelinks,
    pdftitle={Đặc tả và phân tích yêu cầu chức năng phần mềm TMS},
    pdfauthor={Nhóm thực hiện}
}

% =============================================================================
% 5. DINH DANG DE MUC
% =============================================================================
\setlength{\parindent}{0pt}
\setlength{\parskip}{6pt}
\setcounter{secnumdepth}{3}

\renewcommand{\thesection}{\Roman{section}}
\renewcommand{\thesubsection}{\arabic{subsection}}
\renewcommand{\thesubsubsection}{\thesubsection.\arabic{subsubsection}}

\titleformat{\section}{\large\bfseries}{\thesection.}{0.5em}{}
\titleformat{\subsection}{\normalsize\bfseries}{\thesubsection.}{0.5em}{}
\titleformat{\subsubsection}{\normalsize\itshape\bfseries}{\thesubsubsection.}{0.5em}{}

\renewcommand{\cftsecaftersnum}{.}
\renewcommand{\cftsubsecaftersnum}{.}
\renewcommand{\cftsubsubsecaftersnum}{.}

% =============================================================================
% 6. CUSTOM COMMANDS
% =============================================================================
\newcolumntype{C}[1]{>{\centering\arraybackslash}p{#1}}
\newcolumntype{L}[1]{>{\raggedright\arraybackslash}p{#1}}
\newcolumntype{Y}{>{\raggedright\arraybackslash}X}
\newcommand{\code}[1]{{\ttfamily\seqsplit{#1}}}
\newcommand{\tablesetup}{\small\setlength{\tabcolsep}{3pt}\renewcommand{\arraystretch}{1.2}}
\AtBeginEnvironment{longtable}{\small\setlength{\parskip}{0pt}\setlength{\tabcolsep}{3pt}\renewcommand{\arraystretch}{1.15}}

\begin{document}

\setcounter{section}{2}
\section{Đặc tả và phân tích yêu cầu chức năng phần mềm}

Chương này đặc tả các yêu cầu của phần mềm Tutor Management System (TMS) dựa
trên phạm vi đã xác định ở các chương trước. Nội dung được trình bày theo mẫu
báo cáo tham khảo: yêu cầu chức năng nghiệp vụ, biểu mẫu và quy định, bảng
trách nhiệm, yêu cầu phi chức năng, yêu cầu hệ thống, yêu cầu công nghệ và phân
tích khả thi.

\subsection{Yêu cầu chức năng}

\subsubsection{Yêu cầu chức năng nghiệp vụ}

Bảng yêu cầu nghiệp vụ của hệ thống TMS:

\begin{table}[H]
\centering
\tablesetup
\begin{tabular}{|C{0.8cm}|L{3cm}|L{5.1cm}|C{1.4cm}|C{1.4cm}|L{2.1cm}|}
\hline
\textbf{STT} & \textbf{Nghiệp vụ} & \textbf{Thao tác cụ thể} & \textbf{Biểu mẫu} & \textbf{Quy định} & \textbf{Ghi chú} \\
\hline
1 & Quản lý tài khoản & Đăng nhập hệ thống & BM 1.1 & QĐ 1.1 & Xác thực \\
\hline
2 & Quản lý tài khoản & Đăng ký tài khoản giáo viên & BM 1.2 & QĐ 1.2 & Lưu trữ \\
\hline
3 & Quản lý tài khoản & Cập nhật hồ sơ giáo viên và credential Codeforces & BM 1.3 & QĐ 1.3 & Lưu trữ \\
\hline
4 & Quản trị hệ thống & Quản lý tài khoản giáo viên & BM 1.4 & QĐ 1.4 & Admin \\
\hline
5 & Quản trị hệ thống & Cấu hình bot Discord & BM 1.5 & QĐ 1.5 & Tích hợp \\
\hline
6 & Quản lý học sinh & Thêm/cập nhật/tra cứu học sinh & BM 2.1 & QĐ 2.1 & Lưu trữ \\
\hline
7 & Quản lý ghi danh & Ghi danh học sinh vào lớp & BM 2.2 & QĐ 2.2 & Lưu trữ \\
\hline
8 & Quản lý ghi danh & Chuyển lớp, rút khỏi lớp, khôi phục học sinh & BM 2.3 & QĐ 2.3 & Cập nhật trạng thái \\
\hline
9 & Quản lý học sinh & Lưu trữ học sinh chờ xử lý & BM 2.4 & QĐ 2.4 & An toàn dữ liệu \\
\hline
10 & Quản lý lớp học & Thêm/cập nhật/lưu trữ lớp học & BM 3.1 & QĐ 3.1 & Lưu trữ \\
\hline
11 & Quản lý lịch học & Thiết lập lịch học cố định của lớp & BM 3.2 & QĐ 3.2 & Lưu trữ \\
\hline
\end{tabular}
\end{table}

\begin{table}[H]
\centering
\tablesetup
\begin{tabular}{|C{0.8cm}|L{3cm}|L{5.1cm}|C{1.4cm}|C{1.4cm}|L{2.1cm}|}
\hline
\textbf{STT} & \textbf{Nghiệp vụ} & \textbf{Thao tác cụ thể} & \textbf{Biểu mẫu} & \textbf{Quy định} & \textbf{Ghi chú} \\
\hline
12 & Quản lý buổi học & Tạo buổi học thủ công, hủy buổi học, theo dõi trạng thái & BM 3.3 & QĐ 3.3 & Lưu trữ \\
\hline
13 & Quản lý điểm danh & Ghi nhận và cập nhật điểm danh buổi học & BM 3.4 & QĐ 3.4 & Tính phí \\
\hline
14 & Quản lý tài chính & Ghi nhận giao dịch của học sinh & BM 4.1 & QĐ 4.1 & Lưu trữ \\
\hline
15 & Quản lý tài chính & Cập nhật trạng thái khoản phí & BM 4.2 & QĐ 4.2 & Đối soát \\
\hline
16 & Báo cáo tài chính & Xem số dư, tổng quan tài chính và báo cáo thu nhập & BM 4.3 & QĐ 4.3 & Kết xuất \\
\hline
17 & Tích hợp Discord & Liên kết Discord của giáo viên/học sinh & BM 5.1 & QĐ 5.1 & OAuth \\
\hline
18 & Tích hợp Discord & Thiết lập Discord guild/channel cho lớp & BM 5.2 & QĐ 5.2 & Lưu trữ \\
\hline
19 & Tích hợp Discord & Gửi thông báo cho học sinh hoặc lớp học & BM 5.3 & QĐ 5.3 & Gửi tin \\
\hline
20 & Tích hợp Codeforces & Gắn Gym Codeforces cho lớp học & BM 6.1 & QĐ 6.1 & Đồng bộ \\
\hline
21 & Tích hợp Codeforces & Đồng bộ bài tập và bảng xếp hạng Gym & BM 6.2 & QĐ 6.2 & Theo dõi học tập \\
\hline
22 & Báo cáo học tập & Xem tổng quan, hồ sơ học tập và bảng xếp hạng Gym & BM 7.1 & QĐ 7.1 & Tra cứu \\
\hline
\end{tabular}
\end{table}

\subsubsection{Danh sách các biểu mẫu và quy định}

\textbf{BM 1.1: Đăng nhập hệ thống}

\begin{table}[H]
\centering
\tablesetup
\begin{tabularx}{\textwidth}{|C{1cm}|L{4cm}|L{3cm}|Y|}
\hline
\textbf{STT} & \textbf{Trường/Mục} & \textbf{Kiểu dữ liệu} & \textbf{Ghi chú} \\
\hline
1 & Tên đăng nhập & Chuỗi & Tài khoản giáo viên hoặc quản trị viên. \\
\hline
2 & Mật khẩu & Chuỗi & Không hiển thị trực tiếp trên giao diện. \\
\hline
\end{tabularx}
\end{table}

\textbf{QĐ 1.1:} Tên đăng nhập phải tồn tại, mật khẩu phải khớp với mật khẩu đã
băm trong hệ thống, tài khoản phải đang hoạt động. Sau khi đăng nhập thành
công, hệ thống cấp token và xác định đúng vai trò \code{admin} hoặc
\code{teacher}.

\textbf{BM 1.2: Đăng ký tài khoản giáo viên}

\begin{table}[H]
\centering
\tablesetup
\begin{tabularx}{\textwidth}{|C{1cm}|L{4cm}|L{3cm}|Y|}
\hline
\textbf{STT} & \textbf{Trường/Mục} & \textbf{Kiểu dữ liệu} & \textbf{Ghi chú} \\
\hline
1 & Tên đăng nhập & Chuỗi & Dùng để đăng nhập. \\
\hline
2 & Mật khẩu & Chuỗi & Được băm trước khi lưu. \\
\hline
3 & Codeforces handle & Chuỗi & Có thể bỏ trống. \\
\hline
4 & Codeforces API key & Chuỗi & Có thể bỏ trống, lưu dạng bảo vệ. \\
\hline
5 & Codeforces API secret & Chuỗi & Có thể bỏ trống, lưu dạng bảo vệ. \\
\hline
\end{tabularx}
\end{table}

\textbf{QĐ 1.2:} Tên đăng nhập không được trùng; tên đăng nhập và mật khẩu không
được để trống; credential Codeforces nếu nhập phải được lưu an toàn và không
trả nguyên văn ra giao diện.

\textbf{BM 1.3: Cập nhật hồ sơ giáo viên}

\begin{table}[H]
\centering
\tablesetup
\begin{tabularx}{\textwidth}{|C{1cm}|L{4cm}|L{3cm}|Y|}
\hline
\textbf{STT} & \textbf{Trường/Mục} & \textbf{Kiểu dữ liệu} & \textbf{Ghi chú} \\
\hline
1 & Tên đăng nhập mới & Chuỗi & Không bắt buộc. \\
\hline
2 & Mật khẩu mới & Chuỗi & Không bắt buộc. \\
\hline
3 & Codeforces handle & Chuỗi & Cập nhật thông tin luyện tập. \\
\hline
4 & Codeforces API key/secret & Chuỗi & Dùng khi đồng bộ Gym cần quyền truy cập. \\
\hline
\end{tabularx}
\end{table}

\textbf{QĐ 1.3:} Chỉ giáo viên đang đăng nhập được cập nhật hồ sơ của chính
mình; nếu cập nhật mật khẩu thì mật khẩu không được rỗng; dữ liệu nhạy cảm phải
được che khi trả về.

\textbf{BM 1.4: Thông tin tài khoản giáo viên}

\begin{table}[H]
\centering
\tablesetup
\begin{tabularx}{\textwidth}{|C{1cm}|L{4cm}|L{3cm}|Y|}
\hline
\textbf{STT} & \textbf{Trường/Mục} & \textbf{Kiểu dữ liệu} & \textbf{Ghi chú} \\
\hline
1 & Tên đăng nhập & Chuỗi & Duy nhất. \\
\hline
2 & Mật khẩu & Chuỗi & Có thể đặt mới khi tạo hoặc reset. \\
\hline
3 & Trạng thái hoạt động & Boolean & Bật/tắt quyền sử dụng tài khoản. \\
\hline
4 & Codeforces handle/API & Chuỗi & Có thể cập nhật bởi quản trị viên. \\
\hline
\end{tabularx}
\end{table}

\textbf{QĐ 1.4:} Chỉ quản trị viên được quản lý tài khoản giáo viên; không xóa
vật lý tài khoản đã có dữ liệu nghiệp vụ, chỉ thay đổi trạng thái hoạt động khi
cần khóa quyền truy cập.

\textbf{BM 1.5: Cấu hình bot Discord}

\begin{table}[H]
\centering
\tablesetup
\begin{tabularx}{\textwidth}{|C{1cm}|L{4cm}|L{3cm}|Y|}
\hline
\textbf{STT} & \textbf{Trường/Mục} & \textbf{Kiểu dữ liệu} & \textbf{Ghi chú} \\
\hline
1 & Bot token & Chuỗi & Token của Discord bot. \\
\hline
2 & Client ID & Chuỗi & Dùng cho OAuth và link invite. \\
\hline
3 & Client secret & Chuỗi & Dùng cho OAuth callback. \\
\hline
4 & Permissions/scopes & Chuỗi & Quyền cài đặt bot. \\
\hline
\end{tabularx}
\end{table}

\textbf{QĐ 1.5:} Chỉ quản trị viên được cấu hình bot Discord; token và secret
không được hiển thị nguyên văn; hệ thống phải kiểm tra trạng thái bot để cảnh
báo cấu hình sai.

\textbf{BM 2.1: Thông tin học sinh}

\begin{table}[H]
\centering
\tablesetup
\begin{tabularx}{\textwidth}{|C{1cm}|L{4cm}|L{3cm}|Y|}
\hline
\textbf{STT} & \textbf{Trường/Mục} & \textbf{Kiểu dữ liệu} & \textbf{Ghi chú} \\
\hline
1 & Họ tên học sinh & Chuỗi & Bắt buộc. \\
\hline
2 & Lớp hiện tại & Lựa chọn & Lớp thuộc giáo viên đang đăng nhập. \\
\hline
3 & Codeforces handle & Chuỗi & Dùng để đối chiếu kết quả Gym. \\
\hline
4 & Số điện thoại & Chuỗi & Có thể bỏ trống. \\
\hline
5 & Ghi chú & Chuỗi & Có thể bỏ trống. \\
\hline
6 & Ngày ghi danh & Ngày giờ & Thời điểm bắt đầu học. \\
\hline
\end{tabularx}
\end{table}

\textbf{QĐ 2.1:} Họ tên, lớp và ngày ghi danh không được để trống; giáo viên chỉ
được thao tác học sinh thuộc phạm vi của mình; học sinh có trạng thái
\code{active}, \code{pending\_archive} hoặc \code{archived}.

\textbf{BM 2.2: Ghi danh học sinh}

\begin{table}[H]
\centering
\tablesetup
\begin{tabularx}{\textwidth}{|C{1cm}|L{4cm}|L{3cm}|Y|}
\hline
\textbf{STT} & \textbf{Trường/Mục} & \textbf{Kiểu dữ liệu} & \textbf{Ghi chú} \\
\hline
1 & Học sinh & Lựa chọn & Học sinh thuộc giáo viên. \\
\hline
2 & Lớp học & Lựa chọn & Lớp đang hoạt động. \\
\hline
3 & Ngày bắt đầu & Ngày giờ & Mốc bắt đầu ghi danh. \\
\hline
\end{tabularx}
\end{table}

\textbf{QĐ 2.2:} Mỗi học sinh chỉ có tối đa một ghi danh đang hoạt động trong
phạm vi của một giáo viên; lớp được chọn phải đang hoạt động.

\textbf{BM 2.3: Chuyển lớp, rút khỏi lớp, khôi phục học sinh}

\begin{table}[H]
\centering
\tablesetup
\begin{tabularx}{\textwidth}{|C{1cm}|L{4cm}|L{3cm}|Y|}
\hline
\textbf{STT} & \textbf{Trường/Mục} & \textbf{Kiểu dữ liệu} & \textbf{Ghi chú} \\
\hline
1 & Học sinh & Lựa chọn & Học sinh cần cập nhật. \\
\hline
2 & Lớp đích & Lựa chọn & Dùng khi chuyển/khôi phục. \\
\hline
3 & Thời điểm hiệu lực & Ngày giờ & Ngày chuyển, rút hoặc khôi phục. \\
\hline
\end{tabularx}
\end{table}

\textbf{QĐ 2.3:} Khi chuyển lớp, hệ thống đóng ghi danh cũ và tạo ghi danh mới;
khi rút khỏi lớp, hệ thống kết thúc ghi danh hiện tại; nếu còn số dư cần thu
hoặc hoàn, học sinh chuyển sang trạng thái chờ lưu trữ.

\textbf{BM 2.4: Lưu trữ học sinh chờ xử lý}

\begin{table}[H]
\centering
\tablesetup
\begin{tabularx}{\textwidth}{|C{1cm}|L{4cm}|L{3cm}|Y|}
\hline
\textbf{STT} & \textbf{Trường/Mục} & \textbf{Kiểu dữ liệu} & \textbf{Ghi chú} \\
\hline
1 & Học sinh & Lựa chọn & Đang ở trạng thái chờ lưu trữ. \\
\hline
2 & Ngày lưu trữ & Ngày giờ & Thời điểm kết thúc theo dõi. \\
\hline
\end{tabularx}
\end{table}

\textbf{QĐ 2.4:} Chỉ lưu trữ học sinh sau khi đã xử lý các khoản cần thu hoặc
cần hoàn; dữ liệu lịch sử học tập, điểm danh và tài chính vẫn được giữ lại.

\textbf{BM 3.1: Thông tin lớp học}

\begin{table}[H]
\centering
\tablesetup
\begin{tabularx}{\textwidth}{|C{1cm}|L{4cm}|L{3cm}|Y|}
\hline
\textbf{STT} & \textbf{Trường/Mục} & \textbf{Kiểu dữ liệu} & \textbf{Ghi chú} \\
\hline
1 & Tên lớp & Chuỗi & Bắt buộc. \\
\hline
2 & Học phí mỗi buổi & Số thập phân & Dùng sinh khoản phí. \\
\hline
3 & Trạng thái lớp & Lựa chọn & \code{active} hoặc \code{archived}. \\
\hline
\end{tabularx}
\end{table}

\textbf{QĐ 3.1:} Tên lớp không được rỗng; học phí mỗi buổi phải không âm; chỉ
giáo viên sở hữu lớp được cập nhật lớp; lớp đã lưu trữ không được dùng để tạo
ghi danh mới.

\textbf{BM 3.2: Lịch học cố định}

\begin{table}[H]
\centering
\tablesetup
\begin{tabularx}{\textwidth}{|C{1cm}|L{4cm}|L{3cm}|Y|}
\hline
\textbf{STT} & \textbf{Trường/Mục} & \textbf{Kiểu dữ liệu} & \textbf{Ghi chú} \\
\hline
1 & Ngày trong tuần & Số & Từ 0 đến 6 hoặc theo quy ước hệ thống. \\
\hline
2 & Giờ bắt đầu & Giờ & Dạng HH:MM. \\
\hline
3 & Giờ kết thúc & Giờ & Phải sau giờ bắt đầu. \\
\hline
\end{tabularx}
\end{table}

\textbf{QĐ 3.2:} Lịch học phải thuộc một lớp đang hoạt động; giờ kết thúc phải
sau giờ bắt đầu; hệ thống lưu nhiều lịch cố định cho một lớp nếu lớp học nhiều
buổi trong tuần.

\textbf{BM 3.3: Buổi học}

\begin{table}[H]
\centering
\tablesetup
\begin{tabularx}{\textwidth}{|C{1cm}|L{4cm}|L{3cm}|Y|}
\hline
\textbf{STT} & \textbf{Trường/Mục} & \textbf{Kiểu dữ liệu} & \textbf{Ghi chú} \\
\hline
1 & Lớp học & Lựa chọn & Lớp sở hữu buổi học. \\
\hline
2 & Thời điểm bắt đầu & Ngày giờ & Lịch học thực tế. \\
\hline
3 & Giờ kết thúc & Giờ & Có thể lấy theo lịch cố định. \\
\hline
4 & Trạng thái & Lựa chọn & \code{scheduled}, \code{in\_progress}, \code{completed}, \code{cancelled}. \\
\hline
\end{tabularx}
\end{table}

\textbf{QĐ 3.3:} Buổi học chỉ thuộc một lớp; buổi học đã hủy không được sinh
khoản phí mới; buổi học hoàn thành là cơ sở để đối chiếu điểm danh và tài chính.

\textbf{BM 3.4: Điểm danh buổi học}

\begin{table}[H]
\centering
\tablesetup
\begin{tabularx}{\textwidth}{|C{1cm}|L{4cm}|L{3cm}|Y|}
\hline
\textbf{STT} & \textbf{Trường/Mục} & \textbf{Kiểu dữ liệu} & \textbf{Ghi chú} \\
\hline
1 & Buổi học & Lựa chọn & Buổi học cần điểm danh. \\
\hline
2 & Học sinh & Lựa chọn & Học sinh đang ghi danh trong lớp. \\
\hline
3 & Trạng thái điểm danh & Lựa chọn & \code{present}, \code{absent\_excused}, \code{absent\_unexcused}. \\
\hline
4 & Nguồn điểm danh & Lựa chọn & \code{manual}, \code{bot}, \code{system}. \\
\hline
\end{tabularx}
\end{table}

\textbf{QĐ 3.4:} Một học sinh chỉ có một bản ghi điểm danh cho một buổi học;
điểm danh có thể được nhập thủ công hoặc lấy từ bot Discord; trạng thái
\code{present} là căn cứ chính để sinh hoặc giữ khoản phí theo học phí của lớp.

\textbf{BM 4.1: Giao dịch tài chính}

\begin{table}[H]
\centering
\tablesetup
\begin{tabularx}{\textwidth}{|C{1cm}|L{4cm}|L{3cm}|Y|}
\hline
\textbf{STT} & \textbf{Trường/Mục} & \textbf{Kiểu dữ liệu} & \textbf{Ghi chú} \\
\hline
1 & Học sinh & Lựa chọn & Học sinh phát sinh giao dịch. \\
\hline
2 & Số tiền & Số thập phân & Lớn hơn 0. \\
\hline
3 & Loại giao dịch & Lựa chọn & \code{payment} hoặc \code{refund}. \\
\hline
4 & Ngày ghi nhận & Ngày giờ & Có thể là thời điểm hiện tại. \\
\hline
5 & Ghi chú & Chuỗi & Có thể bỏ trống. \\
\hline
\end{tabularx}
\end{table}

\textbf{QĐ 4.1:} Số tiền giao dịch phải lớn hơn 0; học sinh phải thuộc giáo viên
đang đăng nhập; giao dịch được dùng để tính số dư cùng với các khoản phí đang
hoạt động.

\textbf{BM 4.2: Khoản phí}

\begin{table}[H]
\centering
\tablesetup
\begin{tabularx}{\textwidth}{|C{1cm}|L{4cm}|L{3cm}|Y|}
\hline
\textbf{STT} & \textbf{Trường/Mục} & \textbf{Kiểu dữ liệu} & \textbf{Ghi chú} \\
\hline
1 & Khoản phí & Lựa chọn & Phát sinh từ buổi học/điểm danh. \\
\hline
2 & Trạng thái & Lựa chọn & \code{active} hoặc \code{cancelled}. \\
\hline
\end{tabularx}
\end{table}

\textbf{QĐ 4.2:} Chỉ giáo viên sở hữu dữ liệu được cập nhật trạng thái khoản
phí; khoản phí bị hủy không được tính vào tổng phí còn hiệu lực của học sinh.

\textbf{BM 4.3: Tham số báo cáo tài chính}

\begin{table}[H]
\centering
\tablesetup
\begin{tabularx}{\textwidth}{|C{1cm}|L{4cm}|L{3cm}|Y|}
\hline
\textbf{STT} & \textbf{Trường/Mục} & \textbf{Kiểu dữ liệu} & \textbf{Ghi chú} \\
\hline
1 & Ngày bắt đầu & Ngày & Có thể bỏ trống nếu xem toàn bộ. \\
\hline
2 & Ngày kết thúc & Ngày & Phải sau hoặc bằng ngày bắt đầu. \\
\hline
3 & Lớp học & Lựa chọn & Lọc theo lớp nếu cần. \\
\hline
4 & Học sinh & Lựa chọn & Lọc theo học sinh nếu cần. \\
\hline
\end{tabularx}
\end{table}

\textbf{QĐ 4.3:} Báo cáo chỉ hiển thị dữ liệu trong phạm vi giáo viên; khoảng
thời gian nếu nhập phải hợp lệ; số dư học sinh được tính từ tổng giao dịch trừ
tổng khoản phí còn hiệu lực.

\textbf{BM 5.1: Liên kết Discord}

\begin{table}[H]
\centering
\tablesetup
\begin{tabularx}{\textwidth}{|C{1cm}|L{4cm}|L{3cm}|Y|}
\hline
\textbf{STT} & \textbf{Trường/Mục} & \textbf{Kiểu dữ liệu} & \textbf{Ghi chú} \\
\hline
1 & Mã OAuth callback & Chuỗi & Do Discord trả về. \\
\hline
2 & State & Chuỗi & Dùng kiểm tra yêu cầu hợp lệ. \\
\hline
3 & Lỗi OAuth & Chuỗi & Có thể có nếu người dùng từ chối. \\
\hline
\end{tabularx}
\end{table}

\textbf{QĐ 5.1:} OAuth callback phải có state hợp lệ; tài khoản Discord được lưu
đúng với giáo viên hoặc học sinh tương ứng; token học sinh được lưu để phục vụ
thông báo và định danh trong Discord guild.

\textbf{BM 5.2: Thiết lập Discord cho lớp}

\begin{table}[H]
\centering
\tablesetup
\begin{tabularx}{\textwidth}{|C{1cm}|L{4cm}|L{3cm}|Y|}
\hline
\textbf{STT} & \textbf{Trường/Mục} & \textbf{Kiểu dữ liệu} & \textbf{Ghi chú} \\
\hline
1 & Discord guild & Lựa chọn & Guild giáo viên có quyền cấu hình. \\
\hline
2 & Kênh thông báo & Lựa chọn & Kênh text, có thể bỏ trống. \\
\hline
3 & Kênh voice điểm danh & Lựa chọn & Kênh voice, có thể bỏ trống. \\
\hline
\end{tabularx}
\end{table}

\textbf{QĐ 5.2:} Guild được chọn phải thuộc danh sách guild đã đồng bộ của giáo
viên; kênh thông báo phải là text channel; kênh điểm danh phải là voice channel.

\textbf{BM 5.3: Gửi thông báo}

\begin{table}[H]
\centering
\tablesetup
\begin{tabularx}{\textwidth}{|C{1cm}|L{4cm}|L{3cm}|Y|}
\hline
\textbf{STT} & \textbf{Trường/Mục} & \textbf{Kiểu dữ liệu} & \textbf{Ghi chú} \\
\hline
1 & Nội dung & Chuỗi & Nội dung tin nhắn. \\
\hline
2 & Danh sách học sinh & Danh sách & Gửi theo học sinh cụ thể. \\
\hline
3 & Lớp học & Lựa chọn & Gửi cho học sinh trong lớp. \\
\hline
4 & Danh sách guild & Danh sách & Gửi vào kênh đã chọn. \\
\hline
\end{tabularx}
\end{table}

\textbf{QĐ 5.3:} Nội dung không được rỗng; nếu gửi theo lớp thì lớp phải thuộc
giáo viên; hệ thống ghi nhận kết quả gửi thành công hoặc thất bại.

\textbf{BM 6.1: Gắn Gym Codeforces}

\begin{table}[H]
\centering
\tablesetup
\begin{tabularx}{\textwidth}{|C{1cm}|L{4cm}|L{3cm}|Y|}
\hline
\textbf{STT} & \textbf{Trường/Mục} & \textbf{Kiểu dữ liệu} & \textbf{Ghi chú} \\
\hline
1 & Lớp học & Lựa chọn & Lớp cần theo dõi Gym. \\
\hline
2 & Gym ID hoặc liên kết Gym & Chuỗi & Mã Gym trên Codeforces. \\
\hline
3 & Chu kỳ đồng bộ & Số & Phút giữa các lần đồng bộ. \\
\hline
\end{tabularx}
\end{table}

\textbf{QĐ 6.1:} Gym ID không được rỗng; lớp phải đang hoạt động; chu kỳ đồng
bộ phải lớn hơn 0 nếu được nhập; giáo viên cần có credential Codeforces khi Gym
yêu cầu quyền truy cập.

\textbf{BM 6.2: Đồng bộ Gym}

\begin{table}[H]
\centering
\tablesetup
\begin{tabularx}{\textwidth}{|C{1cm}|L{4cm}|L{3cm}|Y|}
\hline
\textbf{STT} & \textbf{Trường/Mục} & \textbf{Kiểu dữ liệu} & \textbf{Ghi chú} \\
\hline
1 & Gym & Lựa chọn & Gym thuộc giáo viên/lớp. \\
\hline
2 & Danh sách bài & Bảng & Problem index và tên bài. \\
\hline
3 & Bảng xếp hạng & Bảng & Kết quả theo học sinh và bài. \\
\hline
\end{tabularx}
\end{table}

\textbf{QĐ 6.2:} Dữ liệu đồng bộ phải gắn đúng Gym; kết quả chỉ đối chiếu với
học sinh có Codeforces handle; thời điểm đồng bộ cuối được lưu để phục vụ theo
dõi.

\textbf{BM 7.1: Tham số báo cáo học tập}

\begin{table}[H]
\centering
\tablesetup
\begin{tabularx}{\textwidth}{|C{1cm}|L{4cm}|L{3cm}|Y|}
\hline
\textbf{STT} & \textbf{Trường/Mục} & \textbf{Kiểu dữ liệu} & \textbf{Ghi chú} \\
\hline
1 & Lớp học & Lựa chọn & Lọc tổng quan theo lớp. \\
\hline
2 & Học sinh & Lựa chọn & Xem hồ sơ học tập. \\
\hline
3 & Gym & Lựa chọn & Xem bảng xếp hạng. \\
\hline
\end{tabularx}
\end{table}

\textbf{QĐ 7.1:} Báo cáo chỉ lấy dữ liệu thuộc giáo viên đang đăng nhập; bảng
xếp hạng Gym hiển thị theo học sinh đang ghi danh trong lớp và danh sách bài đã
đồng bộ.

\subsubsection{Bảng trách nhiệm yêu cầu nghiệp vụ}

\begin{table}[H]
\centering
\tablesetup
\begin{tabular}{|C{0.8cm}|L{3.2cm}|L{4.2cm}|L{6.1cm}|L{1.6cm}|}
\hline
\textbf{STT} & \textbf{Tên yêu cầu} & \textbf{Người dùng} & \textbf{Phần mềm} & \textbf{Ghi chú} \\
\hline
1 & Quản lý tài khoản & Cung cấp thông tin đăng nhập, đăng ký hoặc hồ sơ cần cập nhật. & Kiểm tra dữ liệu, xác thực, băm mật khẩu, cấp token và lưu hồ sơ/credential hợp lệ. & Tự động \\
\hline
2 & Quản trị hệ thống & Quản trị viên cung cấp tài khoản giáo viên hoặc cấu hình bot Discord. & Kiểm tra quyền admin, lưu cấu hình an toàn, che dữ liệu nhạy cảm và kiểm tra trạng thái bot. & Admin \\
\hline
3 & Quản lý học sinh & Giáo viên cung cấp thông tin học sinh, lớp và thời điểm ghi danh. & Kiểm tra phạm vi sở hữu, lưu hồ sơ học sinh, tạo/cập nhật ghi danh và trạng thái học sinh. & Lưu trữ \\
\hline
4 & Quản lý lớp học & Giáo viên cung cấp tên lớp, học phí, lịch học và trạng thái lớp. & Kiểm tra quy định lớp/lịch, lưu lớp, lịch cố định và tạo dữ liệu phục vụ buổi học. & Lưu trữ \\
\hline
5 & Quản lý buổi học và điểm danh & Giáo viên hoặc bot cung cấp trạng thái buổi học và kết quả điểm danh. & Ghi nhận trạng thái buổi học, cập nhật điểm danh, tạo hoặc điều chỉnh khoản phí liên quan. & Tính phí \\
\hline
6 & Quản lý tài chính & Giáo viên cung cấp giao dịch, trạng thái khoản phí hoặc tham số báo cáo. & Kiểm tra số tiền, lưu giao dịch, cập nhật khoản phí, tính số dư và tổng hợp báo cáo. & Đối soát \\
\hline
7 & Tích hợp Discord & Giáo viên/học sinh thực hiện OAuth, chọn guild/channel hoặc nhập nội dung tin nhắn. & Xử lý OAuth callback, lưu định danh Discord, thiết lập binding lớp và gửi thông báo. & Tích hợp \\
\hline
8 & Tích hợp Codeforces & Giáo viên cung cấp credential, Gym ID hoặc yêu cầu đồng bộ. & Gắn Gym, đồng bộ bài tập và bảng xếp hạng, đối chiếu Codeforces handle của học sinh. & Đồng bộ \\
\hline
9 & Báo cáo học tập & Giáo viên chọn lớp, học sinh hoặc Gym cần xem. & Tổng hợp dữ liệu học sinh, lớp, điểm danh, tài chính và Gym để hiển thị báo cáo. & Tra cứu \\
\hline
\end{tabular}
\end{table}

\subsection{Yêu cầu phi chức năng}

\subsubsection{Yêu cầu hiệu quả}

Yêu cầu hiệu quả tập trung vào tốc độ phản hồi của các thao tác thường xuyên và
các thao tác tổng hợp dữ liệu.

\begin{table}[H]
\centering
\tablesetup
\begin{tabularx}{\textwidth}{|C{0.8cm}|L{4cm}|L{5.2cm}|Y|}
\hline
\textbf{STT} & \textbf{Nghiệp vụ} & \textbf{Tốc độ xử lý} & \textbf{Ghi chú} \\
\hline
1 & Đăng nhập và xác thực & Phản hồi dưới 2 giây trong điều kiện kết nối bình thường. & Bao gồm kiểm tra mật khẩu và cấp token. \\
\hline
2 & Tra cứu học sinh/lớp/buổi học & Tải danh sách dưới 5 giây với dữ liệu thông thường. & Có bộ lọc theo trạng thái, lớp, thời gian. \\
\hline
3 & Lưu thông tin nghiệp vụ & Tạo/cập nhật học sinh, lớp, lịch, điểm danh, giao dịch dưới 5 giây. & Không tính thời gian mạng ngoài. \\
\hline
4 & Báo cáo tài chính/học tập & Tổng hợp và hiển thị dưới 15 giây. & Áp dụng cho khoảng thời gian lọc hợp lý. \\
\hline
5 & Đồng bộ Discord/Codeforces & Xử lý nền, không chặn thao tác chính quá 5 giây. & Có thể thông báo trạng thái sau. \\
\hline
\end{tabularx}
\end{table}

\begin{table}[H]
\centering
\tablesetup
\begin{tabularx}{\textwidth}{|C{0.8cm}|L{4cm}|L{4.4cm}|Y|}
\hline
\textbf{STT} & \textbf{Nghiệp vụ} & \textbf{Người dùng} & \textbf{Phần mềm} \\
\hline
1 & Đăng nhập & Cung cấp đúng tên đăng nhập và mật khẩu. & Thực hiện xác thực và trả kết quả trong thời gian quy định. \\
\hline
2 & Quản lý dữ liệu & Nhập đủ dữ liệu theo biểu mẫu. & Kiểm tra nhanh, lưu dữ liệu và phản hồi rõ ràng. \\
\hline
3 & Báo cáo & Chọn bộ lọc phù hợp. & Truy vấn, tổng hợp và hiển thị kết quả đúng phạm vi. \\
\hline
\end{tabularx}
\end{table}

\subsubsection{Yêu cầu tiện dụng}

\begin{table}[H]
\centering
\tablesetup
\begin{tabularx}{\textwidth}{|C{0.8cm}|L{3.4cm}|L{4.4cm}|L{4.4cm}|Y|}
\hline
\textbf{STT} & \textbf{Nghiệp vụ} & \textbf{Mức độ dễ học} & \textbf{Mức độ dễ sử dụng} & \textbf{Ghi chú} \\
\hline
1 & Đăng nhập, hồ sơ & 5 phút hướng dẫn. & Thao tác rõ ràng, lỗi nhập liệu được chỉ ra cụ thể. & Ít trường nhập. \\
\hline
2 & Quản lý học sinh/lớp & 15 phút hướng dẫn. & Danh sách có tìm kiếm, lọc và trạng thái dễ nhận biết. & Dùng thường xuyên. \\
\hline
3 & Điểm danh buổi học & 10 phút hướng dẫn. & Cập nhật nhanh từng học sinh trong một buổi học. & Hạn chế nhầm lớp. \\
\hline
4 & Tài chính & 15 phút hướng dẫn. & Số dư, khoản phí và giao dịch hiển thị thống nhất. & Cần độ chính xác cao. \\
\hline
5 & Discord/Codeforces & 20 phút hướng dẫn. & Có trạng thái liên kết, đồng bộ và lỗi tích hợp. & Phụ thuộc hệ thống ngoài. \\
\hline
6 & Báo cáo & 5 phút hướng dẫn. & Bộ lọc đơn giản, số liệu tổng hợp dễ đọc. & Phục vụ quyết định. \\
\hline
\end{tabularx}
\end{table}

\begin{table}[H]
\centering
\tablesetup
\begin{tabularx}{\textwidth}{|C{0.8cm}|L{3.6cm}|L{4.4cm}|Y|}
\hline
\textbf{STT} & \textbf{Nghiệp vụ} & \textbf{Người dùng} & \textbf{Phần mềm} \\
\hline
1 & Nhập liệu nghiệp vụ & Đọc hướng dẫn và nhập theo biểu mẫu. & Hiển thị nhãn trường, trạng thái bắt buộc và thông báo lỗi dễ hiểu. \\
\hline
2 & Tra cứu và báo cáo & Chọn bộ lọc cần xem. & Cung cấp tìm kiếm, lọc, phân trang và trạng thái trống phù hợp. \\
\hline
3 & Tích hợp ngoài & Thực hiện OAuth hoặc nhập credential theo hướng dẫn. & Điều hướng đến Discord/Codeforces và hiển thị kết quả liên kết rõ ràng. \\
\hline
\end{tabularx}
\end{table}

\subsubsection{Yêu cầu tương thích}

\begin{table}[H]
\centering
\tablesetup
\begin{tabularx}{\textwidth}{|C{0.8cm}|L{4cm}|L{5cm}|Y|}
\hline
\textbf{STT} & \textbf{Nghiệp vụ} & \textbf{Đối tượng liên quan} & \textbf{Ghi chú} \\
\hline
1 & Vận hành giao diện web & Trình duyệt Chrome, Edge, Firefox phiên bản mới. & Giáo viên và quản trị viên sử dụng trên máy tính. \\
\hline
2 & API backend & Frontend TMS và các worker nền. & Giao tiếp qua HTTP nội bộ của ứng dụng. \\
\hline
3 & Cơ sở dữ liệu & Microsoft SQL Server/Azure SQL. & Dùng lưu dữ liệu nghiệp vụ tập trung. \\
\hline
4 & Discord & Discord OAuth, bot API, guild, channel và voice state. & Phụ thuộc token và quyền bot. \\
\hline
5 & Codeforces & Codeforces Gym và API liên quan. & Phụ thuộc handle/credential của giáo viên. \\
\hline
\end{tabularx}
\end{table}

\subsubsection{Yêu cầu tiến hoá}

\begin{table}[H]
\centering
\tablesetup
\begin{tabularx}{\textwidth}{|C{0.8cm}|L{4cm}|L{5.2cm}|Y|}
\hline
\textbf{STT} & \textbf{Nghiệp vụ} & \textbf{Tham số cần thay đổi} & \textbf{Miền giá trị cần thay đổi} \\
\hline
1 & Phân quyền hệ thống & Vai trò và trạng thái tài khoản. & Admin, teacher, active/inactive. \\
\hline
2 & Lớp học & Tên lớp, học phí mỗi buổi, lịch học. & Bảng lớp và lịch học. \\
\hline
3 & Điểm danh & Trạng thái điểm danh và nguồn điểm danh. & present, absent\_excused, absent\_unexcused; manual, bot, system. \\
\hline
4 & Tài chính & Loại giao dịch, trạng thái khoản phí. & payment, refund; active, cancelled. \\
\hline
5 & Discord & Guild/channel được gắn cho lớp. & Danh sách guild/channel đồng bộ. \\
\hline
6 & Codeforces & Gym ID và chu kỳ đồng bộ. & Danh sách Gym và cấu hình đồng bộ. \\
\hline
\end{tabularx}
\end{table}

\begin{table}[H]
\centering
\tablesetup
\begin{tabularx}{\textwidth}{|C{0.8cm}|L{3.6cm}|L{4.3cm}|Y|}
\hline
\textbf{STT} & \textbf{Nghiệp vụ} & \textbf{Người dùng} & \textbf{Phần mềm} \\
\hline
1 & Thay đổi tài khoản & Cho biết tài khoản và trạng thái cần thay đổi. & Ghi nhận trạng thái mới mà không làm mất dữ liệu lịch sử. \\
\hline
2 & Thay đổi lớp/lịch & Cập nhật tên lớp, học phí hoặc lịch học. & Lưu giá trị mới; dữ liệu lịch sử đã phát sinh không bị sửa sai lệch. \\
\hline
3 & Thay đổi cấu hình tích hợp & Chọn guild/channel, credential hoặc Gym mới. & Cập nhật cấu hình và dùng cho các lần xử lý sau. \\
\hline
\end{tabularx}
\end{table}

\subsection{Yêu cầu hệ thống}

\subsubsection{Yêu cầu bảo mật}

Yêu cầu bảo mật đảm bảo người dùng chỉ truy cập đúng dữ liệu thuộc vai trò và
phạm vi của mình.

\begin{table}[H]
\centering
\tablesetup
\begin{tabular}{|C{0.8cm}|L{3.4cm}|L{5.2cm}|C{1.4cm}|C{1.4cm}|C{1.4cm}|}
\hline
\textbf{STT} & \textbf{Nghiệp vụ} & \textbf{Thao tác cụ thể} & \textbf{Admin} & \textbf{Giáo viên} & \textbf{Học sinh} \\
\hline
1 & Quản lý tài khoản & Đăng nhập hệ thống & \checkmark & \checkmark &  \\
\hline
2 & Quản lý tài khoản & Đăng ký/cập nhật hồ sơ giáo viên & \checkmark & \checkmark &  \\
\hline
3 & Quản trị hệ thống & Quản lý tài khoản giáo viên & \checkmark &  &  \\
\hline
4 & Quản trị hệ thống & Cấu hình bot Discord & \checkmark &  &  \\
\hline
5 & Quản lý học sinh & Thêm/cập nhật/tra cứu học sinh &  & \checkmark &  \\
\hline
6 & Quản lý ghi danh & Chuyển lớp, rút lớp, lưu trữ &  & \checkmark &  \\
\hline
7 & Quản lý lớp học & Tạo/cập nhật/lưu trữ lớp và lịch học &  & \checkmark &  \\
\hline
8 & Quản lý buổi học & Tạo/hủy/cập nhật buổi học &  & \checkmark &  \\
\hline
\end{tabular}
\end{table}

\begin{table}[H]
\centering
\tablesetup
\begin{tabular}{|C{0.8cm}|L{3.4cm}|L{5.2cm}|C{1.4cm}|C{1.4cm}|C{1.4cm}|}
\hline
\textbf{STT} & \textbf{Nghiệp vụ} & \textbf{Thao tác cụ thể} & \textbf{Admin} & \textbf{Giáo viên} & \textbf{Học sinh} \\
\hline
9 & Quản lý điểm danh & Ghi nhận điểm danh &  & \checkmark &  \\
\hline
10 & Quản lý tài chính & Giao dịch, khoản phí, số dư &  & \checkmark &  \\
\hline
11 & Tích hợp Discord & Liên kết Discord giáo viên &  & \checkmark &  \\
\hline
12 & Tích hợp Discord & Uỷ quyền Discord học sinh &  &  & \checkmark \\
\hline
13 & Tích hợp Discord & Thiết lập lớp và gửi thông báo &  & \checkmark &  \\
\hline
14 & Tích hợp Codeforces & Gắn Gym, đồng bộ và xem bảng xếp hạng &  & \checkmark &  \\
\hline
15 & Báo cáo & Xem tổng quan, báo cáo tài chính/học tập &  & \checkmark &  \\
\hline
\end{tabular}
\end{table}

\begin{table}[H]
\centering
\tablesetup
\begin{tabularx}{\textwidth}{|C{0.8cm}|L{3.3cm}|L{4.4cm}|Y|}
\hline
\textbf{STT} & \textbf{Vai trò} & \textbf{Người dùng} & \textbf{Phần mềm} \\
\hline
1 & Quản trị viên & Cung cấp thông tin đăng nhập và cấu hình hệ thống. & Chỉ cho phép truy cập chức năng quản trị, không coi admin là giáo viên sở hữu lớp. \\
\hline
2 & Giáo viên & Cung cấp thông tin nghiệp vụ của lớp/học sinh do mình quản lý. & Kiểm tra \code{teacher\_id} ở mọi thao tác đọc/ghi dữ liệu. \\
\hline
3 & Học sinh & Thực hiện uỷ quyền Discord qua link được cấp. & Chỉ lưu credential Discord cho đúng học sinh trong state hợp lệ. \\
\hline
4 & Hệ thống ngoài & Discord/Codeforces trả dữ liệu qua API. & Xử lý token, lỗi API và không công khai credential nhạy cảm. \\
\hline
\end{tabularx}
\end{table}

\subsubsection{Yêu cầu an toàn}

\begin{table}[H]
\centering
\tablesetup
\begin{tabularx}{\textwidth}{|C{0.8cm}|L{3cm}|L{4.2cm}|Y|}
\hline
\textbf{STT} & \textbf{Yêu cầu} & \textbf{Đối tượng} & \textbf{Ghi chú} \\
\hline
1 & Không xóa vật lý & Tài khoản giáo viên & Chỉ thay đổi trạng thái hoạt động để giữ lịch sử dữ liệu. \\
\hline
2 & Không xóa vật lý & Học sinh & Học sinh được chuyển trạng thái lưu trữ; lịch sử ghi danh, điểm danh, tài chính vẫn giữ. \\
\hline
3 & Không xóa vật lý & Lớp học & Lớp được lưu trữ thay vì xóa khi đã có học sinh/buổi học. \\
\hline
4 & Không sửa sai lệch lịch sử & Buổi học đã hoàn thành & Chỉ cho phép điều chỉnh có kiểm soát; dữ liệu tài chính liên quan phải được cập nhật nhất quán. \\
\hline
5 & Không tính khoản phí đã hủy & FeeRecord & Khoản phí \code{cancelled} không được cộng vào số dư phải thu. \\
\hline
6 & Không lộ bí mật tích hợp & Discord token, client secret, Codeforces API key/secret & Chỉ hiển thị trạng thái có/không có credential. \\
\hline
7 & Kiểm tra phạm vi sở hữu & Dữ liệu học sinh, lớp, tài chính, Gym & Giáo viên không được truy cập dữ liệu của giáo viên khác. \\
\hline
\end{tabularx}
\end{table}

\begin{table}[H]
\centering
\tablesetup
\begin{tabularx}{\textwidth}{|C{0.8cm}|L{3.4cm}|L{4.4cm}|Y|}
\hline
\textbf{STT} & \textbf{Nghiệp vụ} & \textbf{Người dùng} & \textbf{Phần mềm} \\
\hline
1 & Lưu trữ tài khoản/học sinh/lớp & Chọn đối tượng cần ngưng sử dụng. & Chuyển trạng thái, giữ dữ liệu lịch sử và ngăn thao tác phát sinh mới nếu không hợp lệ. \\
\hline
2 & Điều chỉnh tài chính & Chọn khoản phí/giao dịch cần cập nhật. & Kiểm tra quyền, lưu trạng thái mới và tính lại số dư. \\
\hline
3 & Credential tích hợp & Nhập credential mới khi cần. & Mã hóa hoặc che dữ liệu nhạy cảm và không trả nguyên văn ra giao diện. \\
\hline
\end{tabularx}
\end{table}

\subsection{Yêu cầu công nghệ}

\begin{table}[H]
\centering
\tablesetup
\begin{tabularx}{\textwidth}{|C{0.8cm}|L{3.4cm}|L{6cm}|Y|}
\hline
\textbf{STT} & \textbf{Yêu cầu} & \textbf{Mô tả chi tiết} & \textbf{Ghi chú} \\
\hline
1 & Nền tảng web & Giao diện chạy trên trình duyệt cho giáo viên và quản trị viên. & Frontend Vite/React. \\
\hline
2 & Backend API & Cung cấp API cho xác thực, học sinh, lớp học, tài chính, Discord, Codeforces. & Node.js/TypeScript. \\
\hline
3 & Kiến trúc backend & Modular monolith, các module giao tiếp trong tiến trình qua contract/public adapter. & Account, System, Student, Classroom, Finance. \\
\hline
4 & Cơ sở dữ liệu & Lưu trữ tập trung bằng Microsoft SQL Server/Azure SQL. & Có ràng buộc khóa ngoại. \\
\hline
5 & Bảo mật & JWT, băm mật khẩu, phân quyền admin/teacher, che credential. & Bắt buộc. \\
\hline
6 & Tích hợp Discord & OAuth, bot token, guild/channel cache, gửi thông báo và hỗ trợ điểm danh voice. & Phụ thuộc quyền bot. \\
\hline
7 & Tích hợp Codeforces & Lưu credential giáo viên, đồng bộ Gym, problems và standings. & Phụ thuộc API Codeforces. \\
\hline
8 & Dễ bảo trì & Mỗi module có \code{domain/application/infrastructure/presentation} riêng. & Giảm ảnh hưởng chéo khi sửa lỗi. \\
\hline
9 & Dễ mở rộng & Có thể bổ sung báo cáo, loại thông báo hoặc nguồn đồng bộ mới. & Không phá vỡ chức năng hiện có. \\
\hline
\end{tabularx}
\end{table}

\begin{table}[H]
\centering
\tablesetup
\begin{tabularx}{\textwidth}{|C{0.8cm}|L{3.4cm}|L{4.4cm}|Y|}
\hline
\textbf{STT} & \textbf{Nghiệp vụ} & \textbf{Người dùng} & \textbf{Phần mềm} \\
\hline
1 & Vận hành frontend & Sử dụng trình duyệt hiện đại. & Gọi API, hiển thị dữ liệu và kiểm tra nhập liệu cơ bản. \\
\hline
2 & Vận hành backend & Gửi yêu cầu qua giao diện. & Xử lý nghiệp vụ, kiểm tra quyền và trả phản hồi chuẩn. \\
\hline
3 & Vận hành cơ sở dữ liệu & Không thao tác trực tiếp. & Đảm bảo ràng buộc, toàn vẹn và truy vấn đúng phạm vi dữ liệu. \\
\hline
4 & Vận hành tích hợp ngoài & Cung cấp credential hoặc OAuth khi cần. & Gọi API Discord/Codeforces, xử lý lỗi và lưu kết quả đồng bộ. \\
\hline
\end{tabularx}
\end{table}

\subsection{Phân tích yêu cầu}

\textbf{Tính khả thi kỹ thuật:} Hệ thống có tính khả thi cao vì các chức năng
cốt lõi đều là các nghiệp vụ web phổ biến: xác thực, quản lý danh mục nghiệp
vụ, điểm danh, tài chính và báo cáo. Rủi ro kỹ thuật chính nằm ở các tích hợp
ngoài như Discord và Codeforces, vì hệ thống phụ thuộc vào token, quyền bot,
OAuth callback và giới hạn API của từng nền tảng.

\textbf{Tính khả thi nghiệp vụ:} Các nghiệp vụ của TMS có ranh giới rõ ràng:
giáo viên quản lý học sinh, lớp, lịch học, buổi học, điểm danh và tài chính;
quản trị viên quản lý tài khoản giáo viên và bot Discord. Dữ liệu phát sinh
theo chuỗi nghiệp vụ tự nhiên từ ghi danh, buổi học, điểm danh đến khoản phí và
số dư, do đó có thể tin học hóa nhất quán.

\textbf{Tính khả thi về con người:} Người dùng chính là giáo viên nên hệ thống
cần ưu tiên giao diện dễ học, thao tác nhanh và báo lỗi rõ ràng. Các chức năng
tích hợp Discord/Codeforces cần có trạng thái liên kết và trạng thái đồng bộ để
giảm nhầm lẫn khi lỗi đến từ hệ thống ngoài.

\textbf{Rủi ro và hướng xử lý:} Rủi ro mất an toàn dữ liệu được giảm bằng cơ chế
phân quyền theo vai trò và theo \code{teacher\_id}. Rủi ro sai lệch tài chính
được giảm bằng việc không xóa vật lý dữ liệu, giữ lịch sử giao dịch/khoản phí
và tính số dư từ các bản ghi nguồn. Rủi ro tích hợp được giảm bằng cách lưu
trạng thái credential, kiểm tra sức khỏe bot và cho phép đồng bộ lại dữ liệu
Codeforces khi cần.

\end{document}
